
GLUE (Wang et al., 2019a), the benchmark that defined NLP evaluation in 2018, was superseded by its own successor within twelve months. When the top-performing model on GLUE surpassed estimated human performance in early 2019, the community did not rely on an automated system to detect this saturation --- it convened, observed that marginal progress had stalled, and collectively decided to create SuperGLUE (Wang et al., 2019b). SuperGLUE itself followed the same arc, prompting BIG-bench \citep{srivastava2022bigbench} and MMLU \citep{hendrycks2021mmlu} roughly two years later. In both cases, community response lagged the actual saturation event by months, during which time substantial research effort was invested in improvements that, retrospectively, offered diminishing capability gains on a test set that was already ''solved.''

This pattern is not a failure of scientific judgment --- it is a failure of measurement infrastructure. Benchmark retirement decisions currently depend on unstructured community observation, meaning the field lacks a principled, automated signal for when a benchmark's utility has been exhausted.

The gap becomes sharper when examining the score-over-time trajectory of these benchmarks. GLUE and SuperGLUE scores do not increase linearly or erratically --- they follow a characteristic S-shaped curve: rapid early improvement, a decelerating middle period, and eventual plateau. This is precisely the shape of a \textbf{logistic growth process}, and it arises naturally from the mechanism of benchmark-specific optimization: early models improve by developing genuine natural language understanding; later models increasingly exploit fixed test set patterns, producing rapidly diminishing marginal gains; eventually all exploitable signal is exhausted. The S-curve is not coincidental --- it is the temporal signature of Goodhart's Law \citep{goodhart1975problems} playing out across a public leaderboard.

This signature is already present in public leaderboard data and is quantitatively detectable without collecting any new test data. Prior approaches to benchmark overfitting most notably require constructing an independent new test set to measure the accuracy gap \citep{recht2019imagenet} --- a resource-intensive, domain-specific intervention that cannot scale to hundreds of benchmarks simultaneously. Qualitative analyses document the problem but offer no automated detection criterion.

This work addresses the gap by reframing benchmark saturation as a \textbf{curve-fitting problem}: given a score-over-time timeseries from a public leaderboard, one can formally characterize the growth trajectory with a logistic model, compare it against alternatives via information-theoretic model selection, and extract an automated saturation date.

The key insight is that bounded logistic fitting (constraining K to near-human-parity levels, allowing negative inflection times) produces overwhelming statistical evidence of S-curve dynamics on both GLUE and SuperGLUE ($\Delta AIC \approx -194$ to $-250$ versus linear and $-113$ to $-357$ versus power-law alternatives --- far exceeding the Burnham \& Anderson (2002) threshold of 10 for decisive model preference), and that a dual-criterion saturation detector applied to the fitted model recovers the community-recognized saturation events within three to five months.

\textbf{Contributions}:

\begin{enumerate}
\item \textbf{A complete automated benchmark saturation detection pipeline} that takes public leaderboard score-over-time data and returns a statistically justified saturation date, requiring no new data collection, human annotation, or domain-specific heuristics.
\end{enumerate}

\begin{enumerate}
\item \textbf{Formal statistical characterization of benchmark score dynamics} demonstrating via AIC model comparison that leaderboard score-over-time trajectories are formally best described by a logistic growth model --- providing, to the authors' knowledge, the first information-theoretic validation of the qualitatively-discussed ''benchmark saturation'' concept.
\end{enumerate}

\begin{enumerate}
\item \textbf{A dual-criterion saturation detector with validated precision}: the detector (top-3 mean $\geq$ 0.99 $\times$ K AND monthly gain $\leq$ 5\% of peak) achieves 3-month accuracy on GLUE and 5-month accuracy on SuperGLUE against community-recognized ground truth. An empirically optimized variant (0.95 $\times$ K threshold) achieves 1-month accuracy on both.
\end{enumerate}

\begin{enumerate}
\item \textbf{Pre-launch inflection discovery}: the logistic inflection point for both GLUE ($t_0 = -6.77$ months) and SuperGLUE ($t_0 = -2.85$ months) predates benchmark launch, indicating that BERT-era pretraining \citep{devlin2019bert} had already saturated benchmark-exploitable NLU signal before systematic leaderboard tracking began.
\end{enumerate}

\begin{enumerate}
\item \textbf{Boundary condition evidence}: the pipeline generalizes to benchmarks with as few as 30 leaderboard entries (mean $R^2$ = 0.913 on 20 synthetic benchmarks), relaxing the initially assumed $\geq$ 50 entry requirement, provided bounded logistic fitting is used.
\end{enumerate}

One negative result is also reported: prospective forecasting from 6-month-early truncated data fails, as the truncated timeseries (approximately 14 months) is insufficient for logistic plateau detection. This scopes the method to retrospective saturation analysis and motivates future work on prospective monitoring with longer lookback windows.

